\subsection{Escolha do empreendimento}

\subsubsection{Público alvo}

Empresas que possuam um corpo de colaboradores de pelo menos dois membros que 
precisem tomar decisões onerosas, críticas e urgentes, como, por
exemplo:

\begin{itemize}
        \item a compra ou venda de um ativo com forte oscilação de preço,
        \item uma movimentação financeira volumosa de urgência para impedir 
            gastos maiores, 
        \item o acionamento de um sistema de segurança que implique em custos, 
        \item o desligamento de um servidor devido a um ataque cibernético, 
            etc.
\end{itemize}

\subsubsection{Definição do problema}

Decisões onerosas, críticas e urgentes, podem comprometer grandes 
volumes de recursos, como: pessoas, recursos financeiros, oportunidades, 
credibilidade, etc. 
Por isso, essas decisões não devem ser tomadas por uma única pessoa, 
deve-se garantir que ela tome uma decisão em concenso com um número 
pré-definido de outros membros da equipe.

Por exemplo, suponha que haja indícios da execução de um ataque de sequestro 
de dados (ransomware), em que criminosos implantam um software malicioso na 
rede da corporação e esse software criptografa os dados dos clientes da 
empresa para posteriormente pedirem o resgate desses dados.

Um colaborador, mexendo no seu terminal num plantão de fim de semana, com 
a maior parte da equipe de folga, não conseguiu acessar várias partes 
do sistema e verificou que alguns arquivos estavam corrompidos. 

Esse colaborador ficaria num dilema, entre duas decisões críticas:
\begin{enumerate}
    \item reunir a equipe urgentemente para que fosse constatado o ataque e 
        fosse tomada a decisão em conjunto, mas o tempo de tomada de decisão 
        poderia potencializar e muito os danos provocados, ou
    \item tomar a decisão de desligar o servidor e ser aclamado herói por ter 
        salvo os dados dos clientes, ou sofrer as consequências de ter 
        tomado uma decisão crítica à revelia e ter deixado o sistema offline.
\end{enumerate}

\subsubsection{Proposta de solução}

No exemplo anterior, certamente, haveria a  necessidade de ferramentas que 
possibilitassem ao identificador da ameaça as seguintes ações:
\begin{itemize}
    \item tocar um alarme para toda a equipe técnica e de negócios,
    \item abrir um chat de discussão de urgência entre os que responderam 
        ao alarme, 
    \item descrever rapidamente a suspeita do ataque,
    \item tomar as providências após o consenso de, digamos, pelo menos três 
        membros da equipe e
    \item gerar um documento com provas digitais do acionamento do alarme e 
        com as assinaturas digitais de todos os envolvidos na tomada de 
        decisão.
\end{itemize}


\subsubsection{Definição do empreendimento}

Uma empresa de soluções para a comunicação estratégica, crítica e de urgência, 
focada em negócios em que a tomada de decisão não pode jamais ser 
responsabilidade de uma única pessoa.
